\begin{figure}[!htbp]
		\centering
		\def\ax{0.5}\def\ay{0.30}
		\providecommand{\xyz}[3]{{#2 + (#1)*\ax},{#3 + (#1)*\ay}}
		\providecommand{\rightangle}[3]{%
			\coordinate (ra1) at ($(#2)!2.4mm!(#1)$);
			\coordinate (ra3) at ($(#2)!2.4mm!(#3)$);
			\draw[line width=0.5pt] (ra1) -- ($(ra1)+(ra3)-(#2)$) -- (ra3);}
		\begin{tikzpicture}[>=Stealth, line width=0.8pt, scale=1.15]
			\draw[gray!60] (\xyz{-1.7}{-1.2}{0}) -- (\xyz{-1.7}{3.90}{0})-- (\xyz{4.7}{3.90}{0}) -- (\xyz{4.7}{-1.2}{0}) -- cycle;
			\node[gray!75] at (\xyz{-1.2}{-1.0}{0}) {$\pi$};
			
			\coordinate (O)   at (\xyz{0}{0}{0});
			\coordinate (Of)  at (\xyz{1.2}{0.3}{0});      % O + v   (base comune)
			\coordinate (Pk)  at (\xyz{0.5}{2.0}{0});      % P            (lambda = 0)
			\coordinate (Ff)  at (\xyz{1.7}{2.3}{0});      % P + v
			\coordinate (Pp)  at (\xyz{2.54}{2.51}{0});    % P'           (lambda = 1.7)
			\coordinate (Ffp) at (\xyz{3.74}{2.81}{0});    % P' + v
			\coordinate (La)  at (\xyz{0.14}{1.91}{0});    % retta d'azione, estremo a
			\coordinate (Lb)  at (\xyz{4.04}{2.89}{0});    % estremo b
			\coordinate (LOa) at (\xyz{-0.48}{-0.12}{0});  % retta per O // v, estremo a
			\coordinate (LOb) at (\xyz{1.92}{0.48}{0});    % estremo b
			\coordinate (MuB) at (\xyz{0}{0}{1.25});
			\coordinate (MuC) at (\xyz{0}{0}{1.47});
			
			% --- le due rette parallele della striscia (tratteggiate, dietro) ---
			%		\draw[gray!55, thin, dashed] (La) -- (Lb);    % retta d'azione di v
			%		\draw[gray!55, thin, dashed] (LOa) -- (LOb);  % retta per O parallela a v
			\draw[gray!55, thin, dashed] ($(O)!-0.6!(Of)$) -- ($(Of)!-1.0!(O)$);
			
			% --- i due parallelogrammi (entrambi grigi): base comune O-Of,
			%     lato opposto sulla stessa retta d'azione => stessa area ---
			\filldraw[gray!9, draw=gray!55, thin] (O) -- (Of) -- (Ffp) -- (Pp) -- cycle; % su OP'
			\filldraw[gray!9, draw=gray!55, thin] (O) -- (Of) -- (Ff)  -- (Pk) -- cycle; % su OP
			
			
			% --- OP e OP' ---
			\draw[->, thin] (O) -- (Pk);
			\node[font=\small, below right] at ($(O)!0.55!(Pk)$) {$\overrightarrow{OP}$};
			\draw[->, thin] (O) -- (Pp);
			\node[font=\small, above left=-4pt] at ($(O)!0.70!(Pp)$) {$\overrightarrow{OP'}$};
			
			% --- v applicato in P e in P' ---
			\draw[->, thick] (Pk) -- (Ff)  node[below right=-3pt] {$\vb$};
			\draw[->, thick] (Pp) -- (Ffp) node[right=-2pt] {$\vb$};
			
			% --- momento mu(O) ---
			\draw[->>, thick] (O) -- (MuB) node[left=1pt, yshift=-10pt] {$\bm{\upmu}(O)$};
			
			\begin{scope}[gray!55]
				\rightangle{MuB}{O}{Of}
				\rightangle{MuB}{O}{Pk}
			\end{scope}
			
			% --- punti ---
			\fill (O)  circle (1.5pt) node[left=2pt] {$O$};
			\fill (Pk) circle (1.5pt) node[below=2pt] {$P$};
			\fill (Pp) circle (1.5pt) node[below right=-2pt] {$P'$};
		\end{tikzpicture}
		\caption{Invarianza del momento polare per traslazione del punto
			di applicazione lungo la linea d'azione.}
		\label{f:invarianza_retta_azione}
	\end{figure}
